\subsection{域上的矩阵}

域 K 上具有 m 行 n 列的有限矩阵与右向量空间 \(E = K_{d}^{n}\) 到右向量空间 \(K_{d}^{m}\) 的线性映射一一对应，这里这些映射的矩阵取关于 E 和 F 的典范基。根据定义，这样的矩阵 X 的秩是与之对应的线性映射 \(u: E \to F\) 的秩；由于这个数按定义是 F 的子空间 \(u(E)\) 的维数，因此（将 X 的列与 u 下 E 的典范基的像等同起来）给出如下定义是等价的：

定义 7. 给定域 K 上 m 行 n 列的矩阵 X，由 X 的 n 个列生成的 \(K_{d}^{m}\) 的子空间的维数称为 X 关于 K 的秩，记作 \(\operatorname{rg}(X)\) 。

也可以说，X 的秩是 X 的线性无关列的最大数目（作为 \(K_{d}^{m}\) 中的元素）。显然 \(\operatorname{rg}(X) \leqslant \inf(m, n)\) ；对 X 的每个子矩阵 Y，有 \(\operatorname{rg}(Y) \leqslant \operatorname{rg}(X)\) 。

若 E 和 F 是 K 上的两个有限维向量空间，且 u 是从 E 到 F 的一个线性映射，则矩阵 \(M(u)\) 关于任意两组基的秩都等于 u 的秩。

命题9. 若一个具有m行n列的矩阵X的元素属于子域 \(K_{0}\) （域K的子域），则X关于 \(K_{0}\) 的秩等于X关于K的秩。

设 \(F_{0}\) 为右向量 \(K_{0}\) -空间，它由右向量 K-空间 \(E = K_{d}^{m}\) 的典范基生成；按假设，X 的列属于 \(E_{0}\) 。设 \(V_{0}\) （相应地，V）为向量子 \(K_{0}\) -空间，它是 \(F_{0}\) 的（相应地，E 的向量子 K-空间的）由这些列生成的子空间。则 \(V = V_{0} \otimes_{K_{0}} K\) （§ 8，第 2 号，命题 2），因此 \(\dim_{K} V = \dim_{K_{0}} V_{0}\) 。

命题10. 域 K 上矩阵 X 的秩等于其转置 \(^{t}X\) 在反域 \(K^{0}\) 上的秩。

在定义 7 之前引入的记号下，u 的秩等于 \({}^{t}u\) 的秩（§7，第 5 号，命题 10），因此该命题由第 4 号命题 3 得出。

由此可见，X 的秩也可以定义为 X 的线性无关行的最大数目（把它们视为左向量 K-空间 \(K_{s}^{n}\) 的元素）。

域 K 上 n 阶方阵与 \(E = K_{d}^{n}\) 的自同态一一对应，并构成一个环，同构于环

End \(_{\kappa}\) (E)（第 7 号）；对应于 E 的自同构的是可逆方阵。

命题11. 设 \(X\) 是一个 \(n\) 阶方阵（域 \(K\) 上的）。以下性质是等价的：

\begin{enumerate}
\def\labelenumi{(\alph{enumi})}
\item
  \(X\) 在 \(\mathbf{M}_n(\mathbf{K})\) 中可逆。
\item
  \(X\) 在 \(\mathbf{M}_n(\mathbf{K})\) 中右可逆。
\item
  \(X\) 在 \(\mathbf{M}_n(\mathbf{K})\) 中左可逆。
\item
  \(X\) 的秩为 \(n\) 。
\end{enumerate}

这只是 §7 第 4 号命题 9 的推论的翻译。

命题12. 对于含 \(m\) 个方程、 \(n\) 个未知数的线性方程组

\[
\sum_ {j = 1} ^ {n} a _ {i j} x _ {j} = b _ {i} \quad (1 \leqslant i \leqslant m)\tag{54}
\]

在域 K 上至少有一个解，必要且充分的条件是：方程组的矩阵 \(A = (a_{ij})\) 与通过给 A 加边一个第 \((n + 1)\) 列（等于 \((b_{i})\) ）而得到的矩阵 B 为同秩的矩阵。

已经看到（第 4 号），(54) 的解的存在性等价于列 \((b_{i})\) 是 A 的各列的线性组合这一事实，因此该命题由 §7 第 3 号命题 4 的推论 4 得出。

注意，当 m = n 且 A 可逆、即秩为 n（命题 11）时，命题 12 的条件总是满足的。若此时 x 和 b 分别表示单列矩阵 \((x_{i})\) 和 \((b_{i})\) ，则方程组 (54) 等价于 A. x = b，且其唯一解为 \(x = A^{-1}.b\) 。

